\chapter{串联与并联电路}\label{ch:g5:series-parallel-circuits}

一个灯泡坏了，整串彩灯全都熄灭了——这就是单回路的“魔咒”。
可是在你家厨房里，灯泡烧坏了，屋里其他的灯依然亮着。
有人把你家的电路接得更巧妙。
今天我们就来学习这种接法：支路接法。

\section{两种接线方式}

你已经认识第一种接法：\emph{串联}，
所有用电器依次串在同一个回路上，像一串项链上的珠子。下面是第二种接法。

\begin{definition}[并联]\label{def:g5:series-parallel-circuits:parallel}
用电器\emph{并联}\index{并联电路}时，电路分成若干条独立的
\emph{支路}\index{支路}，每条支路只接一个用电器，
各支路先分开、后汇合。道路分叉或汇合的地方
称为
\emph{节点}\index{节点}。每条支路都是从电池的一个接线柱
到另一个接线柱的完整通路。
\end{definition}

\begin{omfigure}
\begin{tikzpicture}[scale=0.8]
  \draw (0,0) to[battery1, l=battery] (0,3.2)
    to[short] (2.2,3.2);
  \draw (2.2,3.2) to[short] (2.2,2.4)
    to[lamp, l=bulb 1] (5.4,2.4) to[short] (5.4,3.2);
  \draw (2.2,3.2) to[short] (2.2,4.0)
    to[lamp, l=bulb 2] (5.4,4.0) to[short] (5.4,3.2);
  \draw (5.4,3.2) to[short] (7.0,3.2) to[short] (7.0,0)
    to[short] (0,0);
  \fill (2.2,3.2) circle (2pt);
  \fill (5.4,3.2) circle (2pt);
  \node[above, font=\small] at (3.6,4.95) {两条支路};
\end{tikzpicture}

{\small 两个灯泡并联：电路在一个节点处分成两路，
每个灯泡各占一条支路，两条支路汇合后，
再回到电池。}
\end{omfigure}

\begin{method}[连接方法]\label{met:g5:series-parallel-circuits:build}
用一节电池、两个灯泡和导线连接：
\begin{enumerate}
  \item 先把第一个灯泡单独接在电池上，连成简单回路——
        第一条支路；
  \item 再把第二个灯泡\emph{直接跨接在}第一个灯泡两端：
        从灯泡 1 的每一侧各引一根导线到灯泡 2 的对应一侧——
        第二条支路；
  \item 用手指沿每条支路走一遍：从 $+$ 出发，
        每条路都应恰好经过一个灯泡，然后到达 $-$。
\end{enumerate}
两个灯泡都会亮，
而且每个灯泡都和单独接一个灯泡时一样亮。
\end{method}

\section{支路的规律}

\begin{proposition}[并联电路的规律]\label{prop:g5:series-parallel-circuits:rules}
每个用电器各占一条支路时：
\begin{enumerate}
  \item \emph{独立性}：一条支路断开，只影响
        这一条支路——拧下一个灯泡，另一个灯泡仍然亮着；
  \item \emph{亮度不变}：每个灯泡都如同单独接在电路中一样亮，
        电池的“推力”不会在
        各条支路之间分享；
  \item \emph{控制位置}：支路中的开关只控制本支路；
        装在干路上、分叉之前的开关控制所有用电器，
        （代价是：电池同时给每条支路供电，
        因此耗电更快。）
\end{enumerate}
\end{proposition}

\begin{example}[拧下灯泡的实验再做一次]\label{ex:g5:series-parallel-circuits:test}
对并联的那对灯泡做与串联时完全相同的事：
拧下灯泡 2。灯泡 1 甚至不会闪一下——它自己的支路、
自己完整的通路都未受影响。把灯泡 2 拧回去，改拧下灯泡 1，
灯泡 2 仍然平稳地亮着。单回路的“魔咒”
在并联电路中被解除了：每个灯泡都只管自己的支路。
\end{example}

\begin{omfigure}
\begin{tikzpicture}[scale=0.8]
  \draw (0,0) to[battery1] (0,3.2)
    to[cspst, l=main switch] (2.6,3.2);
  \draw (2.6,3.2) to[short] (2.6,2.4)
    to[lamp, l=bulb 1] (5.8,2.4) to[short] (5.8,3.2);
  \draw (2.6,3.2) to[short] (2.6,4.0)
    to[spst, l=branch switch] (4.2,4.0)
    to[lamp] (5.8,4.0) to[short] (5.8,3.2);
  \draw (5.8,3.2) to[short] (7.4,3.2) to[short] (7.4,0)
    to[short] (0,0);
  \fill (2.6,3.2) circle (2pt);
  \fill (5.8,3.2) circle (2pt);
\end{tikzpicture}

{\small 控制位置：总开关在分叉之前，同时控制两条支路；
支路开关只控制自己的灯泡——此时它是断开的，
所以灯泡 1 单独发光。}
\end{omfigure}

\begin{example}[你家是一个并联电路]\label{ex:g5:series-parallel-circuits:house}
你家里的每盏灯、每个插座和每件电器，都各占一条支路，
共同组成一个大的并联电路。这就是为什么厨房的灯泡烧坏了，
客厅的灯仍然亮着；为什么每个房间的开关只控制自己的灯，
以及为什么电表旁保险丝盒里的总开关，
可以在检修时一次切断全屋的电路。
现代的彩灯也用上了这种接法：它们的灯泡这样连接，
一个灯泡坏了，整串彩灯不会再全部熄灭。
\end{example}

\begin{remark}[串联与并联对比]\label{rem:g5:series-parallel-circuits:compare}
断一处：串联——全部熄灭；并联——只有一条支路熄灭。灯泡增多时：
串联——全部变暗；并联——全部保持正常亮度，
但电池耗电更快。一个开关控制全部：串联——放在哪里都行；并联——只能放在分叉之前。
两种连接都谈不上“更好”。手电筒用串联很合适（一个灯泡，一个开关），
家中电路则一定是并联。电工看任何电路，第一问总是：
电工看任何电路，第一问总是：
电工看任何电路，第一问总是：\emph{道路在何处
分叉？}
\end{remark}

\section{读懂电路}

\begin{method}[沿着路径走]\label{met:g5:series-parallel-circuits:follow}
要预测电路图中任一灯泡的状态：
\begin{enumerate}
  \item 把手指放在电池的 $+$ 极上，
        沿每条可能的路走到 $-$ 极，
        列出每条路经过的元件；
  \item 若至少有一条未断开的路径经过某灯泡，
        该灯泡就会亮；
  \item 同时位于每一条路径上的元件彼此串联；
        位于不同分叉上的元件彼此并联。
\end{enumerate}
预测之前，先把路径走一遍——
只凭图形的样子去猜，正是电路最容易骗人的地方。
\end{method}

\section{练习}

\begin{exercise}[$\star$]\label{exo:g5:series-parallel-circuits:1}
什么是支路？什么是节点？在两个灯泡并联的电路中，
有几条支路上接有灯泡？
\end{exercise}

\begin{exercise}[$\star$]\label{exo:g5:series-parallel-circuits:2}
写出并联电路的三条规律。
\end{exercise}

\begin{exercise}[$\star$]\label{exo:g5:series-parallel-circuits:3}
两个灯泡并联：灯泡 1 被拧下。灯泡 2 会怎样？
串联时，灯泡 2 的情况如何，
两种情况为何不同？
\end{exercise}

\begin{exercise}[$\star$]\label{exo:g5:series-parallel-circuits:4}
两个相同的灯泡：A 组串联，B 组并联，
电池相同。哪一组中每个灯泡更亮？
哪一组的电池耗电更快？
\end{exercise}

\begin{exercise}[$\star$]\label{exo:g5:series-parallel-circuits:5}
在有两个开关的图中，下列情况下各灯泡怎样？
两个开关都闭合？\ 总开关断开？\ 总开关闭合，支路开关
断开？
\end{exercise}

\begin{exercise}[$\star$]\label{exo:g5:series-parallel-circuits:6}
家庭电路为什么采用并联而非串联？
请举出依赖并联的两个日常便利。
\end{exercise}

\begin{exercise}[$\star$]\label{exo:g5:series-parallel-circuits:7}
家庭电路中，哪个开关能同时控制全部用电器？
这个开关日常叫什么？有什么用途？
\end{exercise}

\begin{exercise}[$\star$]\label{exo:g5:series-parallel-circuits:8}
床头灯及其墙上插座，与天花板灯是串联还是
并联？你能从日常生活中怎样判断？
\end{exercise}

\begin{exercise}[$\star\star$]\label{exo:g5:series-parallel-circuits:9}
用符号画出走廊的电路图：一节电池、一个灯泡，
并且还要实现这一便利——灯泡须能在
走廊\emph{任一}端的开关\dots\ 请认真尝试；若你的
设计始终失败，请说明困难所在。（真正的双控
需要比我们所用开关更巧妙的开关——尝试本身
比答案更有启发。）
\end{exercise}

\begin{exercise}[$\star\star$]\label{exo:g5:series-parallel-circuits:10}
一串现代派对彩灯在一只灯泡损坏后仍然发光，
——但盒上警告：“拆去的灯泡不得超过
三只。”请运用支路规律，解释这一便利，
并（难度更高地）说明此警告的合理原因。（想想
每条剩下的支路从电池得到什么，以及
电池自身的负担。）
\end{exercise}

\begin{exercise}[$\star\star$]\label{exo:g5:series-parallel-circuits:11}
在两个开关的图中，电工想再加一个灯泡，
使其\emph{仅}在支路开关闭合时发光，
其余情况下都不亮——即使总开关闭合。新灯泡应接在哪一段
路径上？
接好后，它与灯泡 2 是串联还是并联？
\end{exercise}
